\uuid{RA5x}
\titre{Produit/DL usuels}
\auteur{hanani ; mzari}
\organisation{exo7}
\licence{CC BY-NC-SA 4.0}
\datecreate{2020-01-20}
\niveau{L1}
\module{Analyse}
\chapitre{Développement limité}
\sousChapitre{Calculs}
\difficulte{1}
\theme{Opérations sur les DL, Produit, DL usuels}
\exo7qcmid{lille-2.developpements-limites.003}

\contenu{
\texte{
Soit $a\in \Rr$ et $f$ une fonction. On dit que $f(x)$ est un $o\left((x-a)^n\right)$ au voisinage de $a$ et on écrit $f(x)=o((x-a)^n)$ si $\displaystyle \lim _{x\to a}\frac{f(x)}{(x-a)^n}=0$. On note par $DL_n(a)f(x)$ le développement limité de $f$ à l'ordre $n$ en $a$. On note aussi $DL_n(a^+)f(x)$ (resp. $DL_n(a^-)f(x)$) le développement limité de $f$ à l'ordre $n$ à droite en $a$ (resp. à gauche en $a$).
}

\begin{question}
\qclassification{125.02}
\qtags{motcle=Produit/DL usuels}
Parmi les assertions suivantes, cocher celles qui sont vraies :
\begin{answers}
\good{$\displaystyle DL_3(0)\frac{1}{1-x}=1+x+x^2+x^3+o(x^3)$.}
\bad{$\displaystyle DL_3(0)\frac{1}{1-x}=1-x+x^2-x^3+o(x^3)$.}
\bad{$\displaystyle DL_3(0)\frac{x}{1-x}=x-x^2+x^3+o(x^3)$.}
\good{$\displaystyle DL_3(0)\frac{1+x}{1-x}=1+2x+2x^2+2x^3+o(x^3)$.}
\end{answers}
\begin{explanations}
On a le DL usuel suivant : $\displaystyle \frac{1}{1-x}=1+x+x^2+x^3+o(x^3)$. On en déduit que
$$\frac{x}{1-x}=x(1+x+x^2+x^3+o(x^3))=x+x^2+x^3+o(x^3)$$
et
$$\frac{1+x}{1-x}=(1+x)(1+x+x^2+x^3+o(x^3))=1+2x+2x^2+2x^3+o(x^3).$$
\end{explanations}
\end{question}
}
